\section{Определение матрицы билинейной формы. Формулировка теоремы о преобразовании матрицы билинейной формы при переходе от одного базиса к другому}

\begin{definition}[Матрица билинейной формы]
\textbf{Матрицей билинейной формы} $f$ в базисе $E = \{e_1, \ldots, e_n\}$ называется матрица
\[
A_f^E = \begin{pmatrix}
f(e_1, e_1) & \cdots & f(e_1, e_n) \\
\vdots & \ddots & \vdots \\
f(e_n, e_1) & \cdots & f(e_n, e_n)
\end{pmatrix} = (a_{ij})_{i,j=1}^{n} = \begin{pmatrix}
a_{11} & \cdots & a_{1n} \\
\vdots & \ddots & \vdots \\
a_{n1} & \cdots & a_{nn}
\end{pmatrix}.
\]
\end{definition}

\begin{theorem}[Преобразование матрицы билинейной формы при смене базиса]
Пусть $f(x, y)$ --- билинейная форма, $E = \{e_1, \ldots, e_n\}$ и $E' = \{e'_1, \ldots, e'_n\}$ --- два базиса в $V$, $A_f^E$ и $A_f^{E'}$ --- матрицы формы в этих базисах, $T$ --- матрица перехода от $E$ к $E'$. Тогда
\[
A_f^{E'} = T^T A_f^E T.
\]
\end{theorem}
